% !TEX root = ../main.tex
\section*{摘要}
\addcontentsline{toc}{section}{摘要}

随着外卖平台即时通信功能的广泛使用，用户与骑手之间的订单沟通在提升配送效率的同时，也带来了\textbf{身份信息暴露、订单隐私泄露、聊天记录篡改以及纠纷追责困难}等安全问题。如何在保护用户隐私的同时保留必要的审计与溯源能力，是一个具有实际应用价值的信息安全问题。

针对上述问题，本作品设计并实现了面向外卖订单通信的隐私认证与可信审计系统 FoodShield。系统以\textbf{“日常匿名、异常可溯源”}为设计原则，在不暴露用户真实身份的前提下完成订单认证与实时通信，并通过密码学手段保障通信日志的可信性。当发生投诉或纠纷时，管理员可在权限约束下基于订单号进行条件溯源，在隐私保护与平台管理之间形成平衡。

本作品的主要成果如下：

\begin{enumerate}[leftmargin=2.5em, itemsep=0pt]
\item 面向外卖订单场景设计了可控匿名通信模型。日常通信中隐藏用户真实身份，异常场景下支持审批式条件溯源，避免了完全匿名不可管理与直接实名过度暴露两类问题。
\item 构建了基于国密算法的订单认证与消息保护机制。通信参与方在不暴露真实身份的情况下证明订单合法性，消息明文在存储阶段经过加密处理，日志篡改可被密码学手段检测。
\item 实现了基于 Merkle Tree 的通信日志可信审计。通过快照比对验证历史通信记录的完整性，为订单纠纷处理提供可验证的电子证据。
\end{enumerate}

\vspace{0.5em}
\noindent\textbf{关键词：} 外卖平台；可控匿名；订单认证；Merkle Tree；条件溯源；可信审计
